\documentclass[a4paper,11pt]{article}

% -----------------------------------------------------------------------
% Packages
% -----------------------------------------------------------------------
\usepackage{amsmath,amssymb,amsthm}
\usepackage{hyperref}
\usepackage{booktabs}
\usepackage{microtype}
\emergencystretch=2em
\usepackage{placeins}
\usepackage{xcolor}
\usepackage{listings}

% -----------------------------------------------------------------------
% Theorem environments
% -----------------------------------------------------------------------
\newtheorem{theorem}{Theorem}[section]
\newtheorem{lemma}[theorem]{Lemma}
\newtheorem{corollary}[theorem]{Corollary}
\newtheorem{definition}[theorem]{Definition}
\newtheorem{proposition}[theorem]{Proposition}
\newtheorem{remark}[theorem]{Remark}

% -----------------------------------------------------------------------
% Title
% -----------------------------------------------------------------------
\title{%
  Structural Admissibility for Verification Systems:\\
  Observation, Warrant Debt, and Factorisation Repair%
}

\author{%
  Duston Moore\\
  \texttt{dhwcmoore@gmail.com}
}

\date{}

% -----------------------------------------------------------------------
\begin{document}
\maketitle

% -----------------------------------------------------------------------
\begin{abstract}
Safety claims are usually made through an observation interface: a monitor, a
log, a trace, a sensor, or an abstraction. Such interfaces do not preserve every
concrete distinction. This matters when a safety predicate distinguishes states
that the interface identifies. We call a predicate admissible relative to an
observation map when it factors through that map. If factorisation fails, the
missing distinction cannot be recovered by post-processing the observation. We
formalise this criterion in Rocq. The development contains 41 Rocq source
files, all compiling to corresponding \texttt{.vo} files, with no admitted
lemmas and eight named axioms. The case studies are deliberately
small. They cover consensus under partition, GPU-style relaxed memory
observations, and physical sensor boundary models. They do not claim to be full
models of those domains. Their purpose is narrower: to show how the same
factorisation failure can appear in discrete, concurrent, and physical settings.
The remedy is also ordinary: weaken the claim until the observation supports it,
or strengthen the observation until the distinction is visible. The contribution
is to state this discipline explicitly and to mechanise the reason that some
apparent noise is a failed observation boundary.  The
development also makes the repair obligation operational: extracted
monitor and certificate structures classify admissibility failures as engineering
actions, while warrant debt is classified by kind, repair pressure, and expected
repair burden.
\end{abstract}

\noindent\textbf{Availability.} The Rocq and OCaml sources are available at
\url{https://github.com/dhwcmoore/structural-admissibility}.

% -----------------------------------------------------------------------
\section{Introduction}
\label{sec:introduction}
% -----------------------------------------------------------------------

Verification rarely sees the system directly. It sees an abstraction, a trace, a
log, a monitor state, a sensor reading, or a reconstructed state estimate. Each
is an observation map. Each preserves some distinctions and collapses others.
This is not an error. It is the normal condition under which verification and
runtime assurance operate.

The risk begins when a safety predicate depends on a distinction that the
observation does not preserve. Then the safety claim is stronger than the
interface through which it is made. In routine engineering this problem is often
handled pragmatically. One restricts the operating envelope, adds a sensor,
changes a monitor, rejects an over-strong claim, or records an extra field in a
log. These are familiar moves. The paper gives them a simple structural test.

The same test can be used prospectively.  Given a proposed claim and an
observation interface, the factorisation criterion asks whether the claim is
supported by what the interface actually carries.  If not, the development
classifies the repair obligation: restrict the claim, reject it, or refine the
observation boundary.

A safety predicate is admissible relative to an observation map exactly when it
factors through that map. Equivalently, states that have the same observation
must agree on the predicate. If two concrete states look the same to the
observation but differ with respect to the safety predicate, the claim is not
supported by that observation. No downstream classifier, dashboard, optimiser,
or monitor can repair the loss. The information has already been collapsed.

This point is modest, but important. Many admissibility failures will look like
noise at the observational level. A sensor appears imprecise. A trace appears
partial. A local execution appears harmless. A monitor appears inconclusive. The
paper argues that this is sometimes the wrong diagnosis. The problem may not be
noise. It may be a failed factorisation.

The paper does not propose a new general-purpose verification algorithm. It
does not claim that all safety reasoning reduces to factorisation. It isolates a
necessary condition on safety claims made through observation interfaces. When
that condition fails, the claim must be restricted, rejected, or the observation
must be refined.

\paragraph{Contributions.}

\begin{itemize}
\item We define structural admissibility as factorisation of a safety predicate
through an observation map.
\item We prove the non-recoverability principle: once an observation collapses a
safety-relevant distinction, post-processing cannot separate it.
\item We mechanise the framework in Rocq, with 41 compiled source files, no
admitted lemmas, and eight named axioms.
\item We give small case studies in consensus, GPU-style memory observations,
and physical sensor models.
\item We distinguish two repair forms: predicate rejection or restriction, and
observational refinement. These are formally parallel but not operationally
symmetric.
\item We expose this repair discipline through extracted engineering-action
records, so that monitor and certificate outputs can be read as explicit repair
obligations rather than as unstructured warnings.
\item We repackage the physical threshold assumptions as instances of a
general non-degenerate-threshold obligation, making clear that the issue is
mathematical parameterisation, not an ad hoc feature of the physics example.
\item We refine warrant debt from a mnemonic into a qualitative typology,
classifying debt by kind, repair pressure, and repair cost band.
\end{itemize}

Section~\ref{sec:setting} states the factorisation criterion.
Section~\ref{sec:warrant} defines warrant debt and non-recoverability.
Section~\ref{sec:recovery} explains restriction and refinement.
Section~\ref{sec:rocq} describes the Rocq development.
Section~\ref{sec:casestudies} reports the case studies.
Sections~\ref{sec:related} through~\ref{sec:conclusion} place the result,
state its limits, and draw the practical lesson.

Theorems are proved about the stated models. Industrial relevance follows only
when an engineering setting instantiates the same observation, predicate,
collapse, and repair structure.

% -----------------------------------------------------------------------
\section{Mathematical Setting}
\label{sec:setting}
% -----------------------------------------------------------------------

Let \(X\) be a concrete state space, \(O\) an observation space, and
\(M : X \to O\) an observation map. Let \(\Phi : X \to \mathsf{Prop}\) be a
safety predicate.

\subsection{Observation-induced equivalence}

The observation map \(M\) induces an equivalence on concrete states:
\[
  x \sim_M y \quad\text{iff}\quad M(x) = M(y).
\]
States that agree under \(M\) are indistinguishable from the perspective of
the observation interface.

\subsection{Admissibility}

\begin{definition}[Admissibility]
\label{def:admissible}
The predicate \(\Phi\) is \emph{admissible} relative to \(M\) if and only if
\[
  M(x) = M(y) \Longrightarrow \bigl(\Phi(x) \leftrightarrow \Phi(y)\bigr).
\]
\end{definition}

Equivalently, \(\Phi\) is constant on the fibres of \(M\). This is a support
condition: it says that the observation contains enough information for the
predicate being asserted through it.

\subsection{Factorisation}

\begin{definition}[Factorisation]
\label{def:factorisation}
The predicate \(\Phi\) \emph{factors through} \(M\) if there exists
\(\widehat{\Phi} : O \to \mathsf{Prop}\) such that
\[
  \forall x : X,\quad \Phi(x) \leftrightarrow \widehat{\Phi}(M(x)).
\]
\end{definition}

\subsection{The main equivalence}

\begin{theorem}[Admissibility iff factorisation]
\label{thm:backbone}
The predicate \(\Phi\) is admissible relative to \(M\) if and only if it
factors through \(M\).
\end{theorem}

\begin{proof}
If \(\Phi\) factors through \(M\), then \(M(x)=M(y)\) immediately gives
\(\Phi(x)\leftrightarrow\Phi(y)\). Conversely, suppose \(\Phi\) is constant on
the fibres of \(M\). Define
\[
  \widehat{\Phi}(o) \;:\!\!\iff\; \exists x:X.\ M(x)=o \wedge \Phi(x).
\]
Then \(\Phi(x)\) implies \(\widehat{\Phi}(M(x))\) by taking \(x\) as witness.
Conversely, if \(\widehat{\Phi}(M(x))\) then for some \(y\), \(M(y)=M(x)\) and
\(\Phi(y)\). Fibre-constancy gives \(\Phi(x)\). In the Rocq development this is
\texttt{admissible\_iff\_constant\_on\_kernel}.
\end{proof}

This equivalence is the formal spine of the paper. The rest of the paper is an
application of it. Each case study asks whether a predicate is constant on the
fibres of the stated observation.

% -----------------------------------------------------------------------
\section{Warrant Debt and Non-Recoverability}
\label{sec:warrant}
% -----------------------------------------------------------------------

\subsection{Warrant debt}

\begin{definition}[Witnessed warrant debt]
\label{def:debt}
An observation map \(M\) has \emph{witnessed warrant debt} with respect to
\(\Phi\) when there exist \(x, y : X\) such that
\[
  M(x) = M(y),\quad \Phi(x),\quad \neg\,\Phi(y).
\]
\end{definition}

A witnessed warrant debt is a concrete certificate of non-admissibility. The
term is meant literally. A safety claim made through \(M\) borrows distinctions
that \(M\) does not provide.

\begin{definition}[Profiled warrant debt]
\label{def:profiled-debt}
A \emph{profiled warrant debt} is a witnessed warrant debt together with a
classification of the collapsed distinction and the corresponding repair
pressure.  The classification does not measure risk numerically.  It records
what kind of information has been lost and what sort of engineering act would
be required to discharge the debt.
\end{definition}

\begin{table}[t]
\centering
\small
\begin{tabular}{p{1.9cm}p{2.5cm}p{4.2cm}p{2.0cm}}
\hline
Debt kind & Collapsed distinction & Typical repair & Repair burden \\
\hline
Quorum debt
  & Global quorum structure
  & Enrich quorum data or restrict to local belief
  & Moderate \\[4pt]
Trace debt
  & Global ordering or visibility
  & Add ordering metadata or restrict to DRF executions
  & Low to moderate \\[4pt]
Boundary debt
  & Boundary or inter-sample behaviour
  & Add sensing, margins, or boundary data
  & High \\[4pt]
Threshold debt
  & Non-degenerate threshold condition
  & State threshold obligation explicitly
  & Low, proof-critical \\[4pt]
Runtime debt
  & Fatigue or rupture-state information
  & Emit review or rupture action
  & Operational \\
\hline
\end{tabular}
\caption{Qualitative warrant-debt profiles.  The table is not a quantitative
risk metric.  It classifies the kind of collapsed distinction and the
corresponding repair pressure.}
\label{tab:warrant-debt-profiles}
\end{table}

The point of the profile is not to price safety numerically.  It prevents a
category error.  A thread-local trace debt and a boundary-sensing debt may
instantiate the same formal pattern, but they do not impose the same repair
cost.  The first may require an extra log field or a restricted execution
class.  The second may require additional instrumentation or a stronger
physical model.

\subsection{Non-recoverability}

\begin{theorem}[Non-recoverability]
\label{thm:no-recovery}
For any map \(f : O \to Y\),
\[
  M(x) = M(y) \Longrightarrow f(M(x)) = f(M(y)).
\]
Consequently, if \(M\) collapses \(x\) and \(y\), no downstream map from
observations can separate them.
\end{theorem}

\begin{proof}
By congruence of equality under \(f\). In the Rocq development this is the
non-recoverability lemma for post-processing maps over observations.
\end{proof}

\begin{corollary}[No post-processing repair]
\label{cor:no-repair}
Suppose \(M(x)=M(y)\), \(\Phi(x)\), and \(\neg\Phi(y)\). Then for every
\(f:O\to Y\), the observation \(f\circ M\) does not make \(\Phi\) admissible.
\end{corollary}

\begin{proof}
Since \(M(x)=M(y)\), congruence gives \(f(M(x))=f(M(y))\). But \(\Phi\)
separates \(x\) and \(y\). Hence \(\Phi\) is not constant on the fibres of
\(f\circ M\).
\end{proof}

No function of the observation can recover a distinction not present in the
observation. Classifiers, dashboards, monitors, and optimisers are such
functions unless they receive additional information.

\subsection{Why this can be mistaken for noise}

At the observational level, an inadmissibility witness may look unremarkable.
Two executions have the same local trace. Two sampled trajectories have the
same visible readings. Two distributed states have the same local view. It is
tempting to describe the disagreement as measurement noise, modelling
imprecision, or a benign abstraction artefact. Sometimes that is correct. But
when the predicate separates two states that the observation identifies, the
problem is structural. The observation is not merely noisy with respect to the
claim. It is insufficient for the claim.

This is the practical reason for naming warrant debt. The name does not add a
new proof principle. It marks a common mistake: treating a failed observation
boundary as an implementation nuisance.

\subsection{Consequence for verification}

If a proof of \(\Phi\) requires a distinction invisible to \(M\), the proof is
not a proof through \(M\). It is a proof through a stronger, implicit
observation. The framework makes this precise. An inadmissible predicate cannot
be made admissible by adding an assumption that bypasses the observation
boundary. The assumption merely moves the missing observation into the proof.

% -----------------------------------------------------------------------
\section{Recovery: Restrict or Refine}
\label{sec:recovery}
% -----------------------------------------------------------------------

At the level of the factorisation criterion, there are two principled repair
forms.

\subsection{Predicate rejection or restriction}

One may replace \(\Phi\) by a predicate \(\Psi\) that is admissible relative to
the existing observation \(M\). This may mean weakening the guarantee,
restricting the state class, strengthening the operating envelope, using robust
margins, or limiting executions to those the observation can support.

Rejection does not always mean weakening the predicate in a crude sense. It
means refusing to assert a claim whose truth is not determined by the available
observation. Restriction is the constructive version of that refusal.

\subsection{Observational refinement}

Alternatively, one may replace \(M\) by a richer observation \(M' : X \to O'\)
together with a projection \(\pi : O' \to O\) such that
\[
  M = \pi \circ M'.
\]
The refined observation \(M'\) preserves the old observation while exposing
additional structure. If \(\Phi\) factors through \(M'\), admissibility has been
recovered by changing the observation boundary rather than the predicate.

\subsection{Formal duality is not operational symmetry}
\label{subsec:not-symmetric}

Restriction and refinement are formally parallel repair forms. They are not the
same engineering act. Restriction keeps the observation map fixed and changes
the claim. Refinement keeps the claim fixed and changes the observation. Both
can restore factorisation. They pay different costs.

Restriction pays in guarantee strength. The resulting claim may be narrower,
conditional, local, robust only within a margin, or valid only under a specified
synchronisation discipline. Refinement pays in observational burden. The system
must expose more structure, log more information, add a sensor, preserve a
trace, or reconstruct a richer state.

This distinction is visible in the case studies. Both repairs are mechanised in
all three domains. They are not the same engineering act in all three domains.
The repair that is natural in one setting may be artificial in another.

\subsection{Everyday safety procedures already use both repairs}
\label{subsec:safety-practice}

The distinction is not exotic. Everyday safety procedures already use both
moves. A restricted operating envelope is predicate restriction: the claim is no
longer that all states are safe, but that states inside a stated envelope are
safe. An added sensor is observational refinement: the claim is kept, but the
interface is strengthened. A lockout condition restricts admissible operation. A
new alarm channel refines observation. A conservative threshold restricts the
claim to a margin. A higher resolution measurement refines the observation.

The paper's point is not that engineers do not know this. They do. The point is
that verification texts and monitor designs can obscure the distinction. Once
the distinction is obscured, an unsupported safety predicate can be smuggled in
as a harmless assumption, a modelling convention, or noise. The factorisation
criterion is a discipline against that mistake.

\subsection{Repair choice as an engineering obligation}
\label{subsec:repair-obligation}

The factorisation failure does not decide the repair automatically.  It
creates an obligation.  The engineer must decide whether the unsupported
distinction should be removed from the claim or added to the observation.

\begin{table}[t]
\centering
\small
\begin{tabular}{p{2.2cm}p{2.5cm}p{2.8cm}p{3.5cm}}
\hline
Repair & Claim preserved? & Observation changed? & Cost paid \\
\hline
Restriction & No & No & Reduced guarantee strength \\[4pt]
Rejection   & No & No & Claim abandoned \\[4pt]
Refinement  & Yes & Yes & Added observation burden \\[4pt]
Explicit carriage & No full discharge & No & Documented limitation \\
\hline
\end{tabular}
\caption{Repair choices after factorisation failure.  Restriction and refinement
are formally parallel repairs, but they are not operationally symmetric.}
\label{tab:repair-choice}
\end{table}

A restricted claim may be the correct engineering response when the stronger
claim is not worth the instrumentation required to support it.  A refined
observation is the correct response when the stronger claim is required and
the missing distinction can be exposed by logging, sensing, reconstruction,
or synchronisation discipline.  Explicit carriage is not a proof repair.  It
is the act of recording that a known admissibility debt remains attached to a
weaker safety case.

\subsection{No invisible observation}

\begin{quote}
\textbf{Principle 4.1} (No invisible observation). When a predicate fails to
factor through the current observation map, a proof must not be repaired by an
assumption that supplies the missing distinction without changing the stated
observation boundary.
\end{quote}

The proof obligation is discharged only by restricting the claim, rejecting the
claim, or refining the observation. The case studies in
Section~\ref{sec:casestudies} illustrate this discipline.

% -----------------------------------------------------------------------
\section{Rocq Development}
\label{sec:rocq}
% -----------------------------------------------------------------------

\subsection{Library structure}

The mechanisation is organised as follows:

\begin{description}
\item[\texttt{Foundation/}] Relations, orders, lattices, quotients, topology,
mereotopology, and metric resolution.
\item[\texttt{Admissibility/}] Observation maps, the factorisation criterion,
kernel pairs, refinement, collapse witnesses, warrant debt, profiled warrant
debt, and horizon debt.
\item[\texttt{Runtime/}] Monitor states, structural fatigue, metastable closure,
rupture certificates, recovery, and runtime soundness.
\item[\texttt{Automation/}] Custom tactics for admissibility, refinement,
collapse, and certificate reasoning.
\item[\texttt{CaseStudies/Consensus/}] Network model, partition collapse,
consensus admissibility.
\item[\texttt{CaseStudies/GPU/}] Memory model, GPU admissibility.
\item[\texttt{CaseStudies/PhysicalSystems/}] Boundary conditions,
non-degenerate threshold obligations, thermal model, sensor collapse, PDE
admissibility.
\item[\texttt{Extraction/}] Extractable types, monitor extraction, certificate
extraction, engineering-action extraction, and extraction correctness.
\end{description}

\subsection{Formal status}

The accompanying Rocq development consists of 41 source files, all of which
compile to \texttt{.vo} files. It contains no admitted lemmas. The remaining
logical strength of the development is confined to eight named axioms. They are
not hidden. They are part of the mathematical specification.

\subsection{Axioms and their role}
\label{subsec:axioms-role}

The eight axioms are background principles, not incomplete proofs. Three concern
primitive connection in mereotopology:

\begin{description}
\item[\texttt{C\_refl}] Every region is connected to itself.
\item[\texttt{C\_sym}] Connection is symmetric.
\item[\texttt{ntpp\_irreflexive}] No region is a non-tangential proper part of
itself.
\end{description}

The axiom \texttt{ntpp\_irreflexive} is an RCC-style mereotopological
assumption. It is not derivable from \texttt{C\_refl} and \texttt{C\_sym}
alone. It is therefore stated explicitly. The derived result
\texttt{ntpp\_antisym} is proved from it.

The remaining five axioms are application-level constraints. Their exact Rocq
names and roles are given in Table~\ref{tab:axiom-roles}. The following
paragraphs explain how each axiom interacts with the actual proof obligations
it discharges.

\begin{table}[h]
\centering
\small
\begin{tabular}{p{3.8cm}p{4.0cm}p{3.2cm}}
\toprule
Rocq axiom name & Role in the model & Why it is explicit \\
\midrule
\texttt{C\_refl}
  & Every region connects to itself
  & Base mereotopology \\[3pt]
\texttt{C\_sym}
  & Connection is symmetric
  & Base mereotopology \\[3pt]
\texttt{ntpp\_irreflexive}
  & No region is an NTPP of itself
  & Boundary discipline \\[3pt]
\texttt{boundary\_obs\_\allowbreak{}respects\_\allowbreak{}signature}
  & Same boundary signature implies same observation
  & Sensor interface coherence \\[3pt]
\texttt{max\_temp\_positive}
  & Thermal safety threshold is positive
  & Non-degenerate witness \\[3pt]
\texttt{max\_pressure\_positive}
  & Pressure safety threshold is positive
  & Non-degenerate witness \\[3pt]
\texttt{max\_flux\_positive}
  & Flux safety threshold is positive
  & Non-degenerate witness \\[3pt]
\texttt{meta\_review\_\allowbreak{}trigger\_lt\_\allowbreak{}fatigue\_threshold}
  & Review trigger fires before hard rupture threshold
  & Runtime monitor ordering \\
\bottomrule
\end{tabular}
\caption{All eight named axioms with their exact Rocq names and roles.
  Full justification for each appears in \texttt{ASSUMPTIONS.md}.}
\label{tab:axiom-roles}
\end{table}

\paragraph{\texttt{boundary\_obs\_respects\_signature}.}%
This axiom is the contractual interface between the abstract sensor model and
the factorisation failure proof. It states that regions sharing the same
boundary signature produce identical boundary observations. The key use is in
\texttt{boundary\_collapse\_induces\_\allowbreak{}admissibility\_failure}: the proof must
rewrite the boundary observation of one region in terms of the other after
establishing that their signatures agree. Without the axiom, that rewrite is
not licensed. The axiom is not a consequence of the mereotopological base. It
is an explicit requirement imposed on any concrete instantiation of the boundary
sensor interface.

This is the physical analogue of a hidden structural bridge in the discrete
case studies.  In the GPU case, the restriction proof relies on the DRF-SC
bridge for the toy execution model.  In the physical case, the corresponding
burden is boundary coherence and non-degenerate thresholds.  The burdens
differ in content, but they play the same methodological role: they state the
conditions under which the observation map is precise enough for the
factorisation question to be well posed.

\paragraph{\texttt{max\_temp\_positive}, \texttt{max\_pressure\_positive},
\texttt{max\_flux\_positive}.}
These three axioms are instances of one structural obligation: a threshold
used to build a predecessor witness must be non-degenerate.  The proofs
compare a threshold value with its predecessor.  In Rocq's natural-number
arithmetic this comparison is meaningful only when the threshold is strictly
positive.  Otherwise subtraction by one saturates at zero, and the intended
witness degenerates.

The development therefore packages the issue as a non-degenerate-threshold
obligation.  Temperature, pressure, and flux are three instances of the same
mathematical requirement.  The point is not that the physical model needs
three unrelated physical assumptions.  The point is that any threshold-based
safety witness of this form must state the arithmetic condition under which
the witness is non-degenerate.

\paragraph{\texttt{meta\_review\_trigger\_lt\_fatigue\_threshold}.}
This axiom orders the two key parameters of the runtime monitor: the fatigue
level at which the monitor enters \textsc{Meta\_Review}, and the level at which
it enters \textsc{Hard\_Rupture}. The strict inequality ensures that there
exists a fatigue band in which the monitor is under structural review but has
not yet triggered hard rupture. Without this ordering, the
\textsc{Meta\_Review} state could be unreachable: if the hard-rupture threshold
were not above the review trigger, the monitor could cross directly from
\textsc{Closed} to \textsc{Hard\_Rupture}, bypassing the intermediate phase.
The axiom governs the topology of the monitor state machine: it is the
condition that makes metastable closure distinct from hard rupture.

The continuous case does more mathematical work than the discrete examples:
boundary assumptions and ordering constraints have no direct analogue in a
thread-local trace or a partitioned local view. The common claim is that each
domain contains an observation map, a predicate, and a possible failure of
factorisation.

\subsection{Mechanisation scope}

The development mechanises the factorisation criterion and its equivalence with
admissibility, the non-recoverability principle, the runtime monitor theory
including rupture certificates, and the three case-study examples. It does not
claim to be a full mechanisation of Paxos, a production GPU memory model, or a
continuous control-system specification.

Table~\ref{tab:code-claims} maps each prose claim to its Rocq artefact.

\begin{table}[h]
\centering
\small
\begin{tabular}{p{3.5cm}p{4.5cm}p{2.6cm}}
\toprule
Claim & Rocq artefact & Status \\
\midrule
Admissibility iff fibre-constancy
  & \texttt{admissible\_iff\_\allowbreak{}constant\_on\_kernel}
  & proved \\[3pt]
Non-recoverability
  & non-recoverability lemma in \texttt{Collapse.v}
  & proved \\[3pt]
Consensus refinement
  & \texttt{ConsensusAdmissibility.v}
  & proved \\[3pt]
Consensus restriction
  & \texttt{ConsensusRestriction.v}
  & proved \\[3pt]
GPU refinement
  & \texttt{GPUAdmissibility.v}
  & proved \\[3pt]
GPU DRF restriction
  & \texttt{GPURestriction.v}
  & proved, conditional on DRF-SC bridge \\[3pt]
Physical restriction
  & \texttt{PDEAdmissibility.v}
  & proved \\[3pt]
Physical refinement
  & \texttt{PDEAdmissibility.v}
  & proved \\[3pt]
Extraction bridge
  & \texttt{ExtractionCorrectness.v}
  & proved \\[3pt]
Profiled warrant debt
  & \texttt{WarrantDebt.v}
  & proved/\allowbreak{}classified \\[3pt]
Non-degenerate thresholds
  & \texttt{BoundaryConditions.v}, \texttt{ThermalModel.v}, \texttt{PDEAdmissibility.v}
  & parameterised \\[3pt]
Engineering-action extraction
  & \texttt{MonitorExtraction.v}, \texttt{CertificateExtraction.v}
  & extracted \\
\bottomrule
\end{tabular}
\caption{Prose claims and their Rocq artefacts.}
\label{tab:code-claims}
\end{table}

The profiling and action records are not additional safety theorems about the
external world.  They are typed interfaces around the existing factorisation
result.  Their role is to prevent an admissibility failure from being reported
as an unstructured warning.  The reported object records the kind of debt, the
repair pressure, and the corresponding engineering action.

\subsection{Extraction bridge}
\label{subsec:extraction}

The \texttt{Extraction/} component produces two OCaml files from the Rocq
development via the built-in \texttt{Extraction Language OCaml} command:
\texttt{monitor\_extracted.ml} and \texttt{cert\_extracted.ml}.

\texttt{monitor\_extracted.ml} contains the function \texttt{update\_monitor},
which takes an \texttt{ExtractableSnapshot} and an \texttt{ExtractableInput}
and returns the next snapshot. The Rocq type \texttt{ExtractableSnapshot}
records fatigue level, tension, horizon, lag, and numeric status code (0 =
\textsc{Closed}, 1 = \textsc{Meta\_Review}, 2 = \textsc{Hard\_Rupture}). The
extracted file also contains the proofs of two verified properties:
\texttt{update\_monitor\_preserves\_invariants}, which states that the monitor
invariant is maintained across every transition, and
\texttt{update\_monitor\_hard\_rupture\_absorbing}, which states that once the
hard-rupture status is set it is never cleared. The proof terms carry no
operational content in the extracted program; the computational content is the
monitor update function itself. The extracted code is therefore not a proof
checker. It is code produced from terms whose relevant properties have been
proved in Rocq.

\texttt{cert\_extracted.ml} contains \texttt{extracted\_cert\_valid}, a
boolean certificate checker, and \texttt{build\_extractable\_cert}, a
constructor for rupture certificates from explicit witnesses. Within the Rocq model, the soundness result
\texttt{extracted\_certificate\_sound} states that every certificate accepted
by the checker satisfies the certificate predicate: the two predicate values
recorded in the certificate differ, and the fatigue component is positive. It
does not certify the external sensor pipeline, the OCaml compiler, or the
runtime environment.

The correctness connection between the Rocq specification and the extracted
code rests on three proved theorems in \texttt{ExtractionCorrectness.v}.
\texttt{snapshot\_round\_trip} proves that converting a \texttt{MonitorSnapshot}
to the extractable representation and back is an identity. The theorem
\texttt{coq\_extractable\_coq\_equiv} proves that a round-trip through the
extractable types preserves both fatigue level and status code. Together they ensure that the type bridge is faithful: no information is
discarded or corrupted by the representation change.

The extraction layer also exposes an engineering-action interface.  A monitor
status is mapped to an action classification: closed states require no repair
action, meta-review states require review of the observation boundary, and
hard-rupture states require rejection or refinement.  Certificate output is
likewise mapped to a repair obligation when it records a concrete predicate
disagreement over an identical observation.  This is not a dashboard
implementation.  It is the typed core that a dashboard, issue tracker, or
safety-case tool would consume.

\begin{table}[t]
\centering
\small
\begin{tabular}{p{2.5cm}p{3.2cm}p{5.0cm}}
\hline
Extracted condition & Engineering action & Interpretation \\
\hline
Closed
  & No repair action
  & Current observation supports the monitored state \\[4pt]
Meta Review
  & Review observation boundary
  & Structural fatigue is visible before rupture \\[4pt]
Hard Rupture
  & Reject, restrict, or refine
  & The current claim is no longer supportable \\[4pt]
Certificate witness
  & Open admissibility obligation
  & Same observation, different predicate value \\
\hline
\end{tabular}
\caption{Operational reading of extracted monitor and certificate output.}
\label{tab:engineering-action}
\end{table}

The limitation of the extraction layer is the standard one. Rocq's extraction
theorem guarantees that the extracted OCaml has the same computational
behaviour as the Rocq term under the Rocq reduction rules. It does not verify
the OCaml compiler, the runtime system, or the interface code that feeds
sensor data into \texttt{ExtractableInput} records. That boundary is the point
where the formal guarantee ends and engineering practice resumes.

% -----------------------------------------------------------------------
% \section{Case Studies} — title comes from case-studies.tex
\label{sec:casestudies}
% -----------------------------------------------------------------------

% case-studies.tex
% Integrated case-study section for the structural admissibility paper.
% Insert into the main file with:
%   % case-studies.tex
% Integrated case-study section for the structural admissibility paper.
% Insert into the main file with:
%   % case-studies.tex
% Integrated case-study section for the structural admissibility paper.
% Insert into the main file with:
%   \input{paper/case-studies}
% at the point where the case-study section begins.

\section{Case Studies: Refine or Reject}

The development contains three case-study families: consensus under network
partition, GPU-style relaxed memory, and cyber-physical sensor models. These
examples are deliberately drawn from different engineering domains. Their common
structure is the point. In each case, a safety predicate is naturally formulated
over a concrete state space~\(X\), while the available observation map
\(M : X \to O\) exposes only a coarser interface. The case study then asks
whether the predicate factors through that observation.

The answer is negative for the initial coarse observation in each family. The
coarse observation identifies concrete states that disagree on the relevant
safety predicate. The admissibility failure is therefore structural rather than
accidental. Recovery takes one of two forms: replace the predicate by one that
is meaningful at the observational grain, or refine the observation so that the
safety-relevant distinction is no longer collapsed.

\begin{table}[t]
\centering
\caption{Case-study structure: factorisation failure and admissibility recovery.}
\label{tab:case-study-factorisation}
\small
\begin{tabular}{p{1.8cm}p{2.8cm}p{3.5cm}p{2.9cm}}
\hline
Case study & Hidden distinction & Coarse failure & Recovery \\
\hline
Consensus
  & Partition and quorum structure
  & Local observations do not determine global quorum safety
  & Restrict to local quorum belief; or enrich with quorum-safety data \\[4pt]
GPU memory
  & Buffer, cache, and ordering structure
  & Value traces do not determine consistency behaviour
  & Restrict to DRF executions; or enrich with global ordering metadata \\[4pt]
Physical systems
  & Inter-sample and boundary behaviour
  & Sampled values do not determine continuous safety
  & Use robust observational margins; or enrich sensing and boundary data \\
\hline
\end{tabular}
\end{table}

\paragraph{Mechanisation status.}
The three case-study families instantiate the same factorisation pattern.
The consensus case study mechanises both repair paths: the refinement
path restores admissibility for \(\texttt{quorum\_safe}\) by enriching
the observation (\texttt{ConsensusAdmissibility.v}), and the restriction
path yields an admissible local-quorum-belief predicate at the coarse
observational grain (\texttt{ConsensusRestriction.v}). The GPU case study
mechanises both repair paths: refinement by enriched global execution
observation (\texttt{GPUAdmissibility.v}), and restriction to the subtype
of data-race-free executions (\texttt{GPURestriction.v}). The
physical-systems case study mechanises both paths: a restricted predicate
is admissible at the coarse sensor grain, and an enriched sensor observation
restores admissibility for the stronger predicate. The development
mechanises predicate restriction in all three domains and observational
refinement in all three.

% -----------------------------------------------------------------------
\subsection{Consensus Under Partition}

The consensus case study concerns quorum safety under network partition. The
concrete state records enough information to determine whether the relevant
quorum-safety predicate holds. A coarse observation, however, may expose only
local status information. Such an observation does not determine global quorum
safety. Two concrete states may agree on the local observation while disagreeing
on whether a quorum is safely established.

This failure is not a defect in the proof script. It is the phenomenon the
framework is intended to detect. A uniformity principle for reachable alive
nodes fails in the partition model: different nodes may have different
reachable alive counts. The admissibility proof cannot be closed by that
route.

The mechanised consensus case study proves both repair paths.

\textit{Refinement.} The observation record \texttt{ConsensusObs} contains the
field \texttt{co\_quorum\_safe}, computed by \texttt{consensus\_obs} from the
underlying consensus state. The proof of admissibility is then immediate:
equality of observations entails equality of the quorum-safety component by
projection. The safety predicate becomes admissible because the observation
carries the safety-relevant distinction directly.

\textit{Restriction.} Rather than asserting global quorum safety, one can
instead assert only what the local observation supports: the node's own quorum
belief, \texttt{local\_quorum\_belief~n}, a projection of \texttt{lv\_quorum}
from the local view record. This predicate is admissible relative to
\texttt{local\_node\_obs~n} because it is a projection of the observation.
The proof requires no axiom. The theorem \texttt{restriction\_is\_genuinely\_weaker}
confirms that local belief and global quorum safety can disagree, so this is a
genuine weakening and not a reformulation.

The theorem \texttt{both\_repairs\_are\_available} states the formal duality
in a single proposition: both \texttt{local\_quorum\_belief~n} and
\texttt{quorum\_safe} are admissible, the former relative to the coarse local
observation, the latter relative to the enriched consensus observation.

% -----------------------------------------------------------------------
\subsection{GPU-Style Memory Observations}

The GPU-style case study concerns the gap between value observations and memory
consistency. A coarse observation may record only the values read and written by
threads. The concrete execution state, however, contains additional ordering and
visibility structure: buffering, cache state, scope, and synchronisation
metadata. A predicate such as sequential consistency of the concurrent execution
is not, in general, determined by the value trace alone.

The admissibility failure has the same form as in the consensus case. Two
executions may have the same observed value trace for a given thread while
disagreeing on whether the execution is sequentially consistent. The observation
map has collapsed a distinction on which the safety predicate depends. The
predicate therefore fails to factor through the thread-local observation.

The GPU case study mechanises both repair paths. Refinement restores
admissibility by enriching the observation with global execution structure.
Restriction keeps the thread-local observation fixed but restricts the domain
to data-race-free executions. On that restricted domain, the
sequential-consistency predicate is stable across observation fibres,
conditional on the DRF-SC bridge stated in \texttt{GPURestriction.v}.

% -----------------------------------------------------------------------
\subsection{Physical Sensing and Boundary Conditions}

The physical-systems case study concerns safety predicates over continuous or
boundary-sensitive states. A controller or monitor rarely observes the full
physical trajectory. It receives sampled, discretised, or otherwise compressed
sensor data. A predicate may require that a variable remain within a safe region
throughout an interval, while the observation records only values at selected
sampling points or coarse boundary summaries.

The resulting admissibility failure is again a failure of factorisation. Two
physical trajectories may induce the same sampled observation while disagreeing
on the safety predicate. One trajectory may remain inside the safe envelope,
while another may violate the envelope between samples or across an unresolved
boundary. The sampled observation is therefore insufficient to determine the
predicate.

This case study mechanises both recovery paths. The coarse sensor observation is
proved inadmissible for the original safety predicate. A restricted predicate
meaningful at the coarse observational grain is then shown to be admissible
relative to that observation. Finally, an enriched sensor observation that adds
the missing physical data restores admissibility for the stronger predicate. The
physical-systems case study therefore demonstrates the full design choice: one
may weaken the guarantee to match the observation, or strengthen the observation
to match the guarantee.

% -----------------------------------------------------------------------
\subsection{Common Structure}

The three case-study families differ in engineering content but not in logical
shape. Consensus hides partition and quorum structure. GPU-style memory hides
ordering and visibility structure. Physical sensing hides inter-sample or
boundary behaviour. In each case, the initial observation map is too coarse for
the safety predicate. The predicate fails to factor because the observation
identifies concrete states that the predicate separates.

This is the common admissibility defect. It is independent of whether the hidden
structure is distributed, concurrent, or physical. The repair is also common:
either choose a predicate that is meaningful at the given observational grain, or
refine the observation until the predicate becomes invariant on observation
fibres. The case studies therefore support the main methodological claim of the
paper. Admissibility is not merely a property of a predicate or a system. It is a
property of the relation between a predicate and an observation interface.

% at the point where the case-study section begins.

\section{Case Studies: Refine or Reject}

The development contains three case-study families: consensus under network
partition, GPU-style relaxed memory, and cyber-physical sensor models. These
examples are deliberately drawn from different engineering domains. Their common
structure is the point. In each case, a safety predicate is naturally formulated
over a concrete state space~\(X\), while the available observation map
\(M : X \to O\) exposes only a coarser interface. The case study then asks
whether the predicate factors through that observation.

The answer is negative for the initial coarse observation in each family. The
coarse observation identifies concrete states that disagree on the relevant
safety predicate. The admissibility failure is therefore structural rather than
accidental. Recovery takes one of two forms: replace the predicate by one that
is meaningful at the observational grain, or refine the observation so that the
safety-relevant distinction is no longer collapsed.

\begin{table}[t]
\centering
\caption{Case-study structure: factorisation failure and admissibility recovery.}
\label{tab:case-study-factorisation}
\small
\begin{tabular}{p{1.8cm}p{2.8cm}p{3.5cm}p{2.9cm}}
\hline
Case study & Hidden distinction & Coarse failure & Recovery \\
\hline
Consensus
  & Partition and quorum structure
  & Local observations do not determine global quorum safety
  & Restrict to local quorum belief; or enrich with quorum-safety data \\[4pt]
GPU memory
  & Buffer, cache, and ordering structure
  & Value traces do not determine consistency behaviour
  & Restrict to DRF executions; or enrich with global ordering metadata \\[4pt]
Physical systems
  & Inter-sample and boundary behaviour
  & Sampled values do not determine continuous safety
  & Use robust observational margins; or enrich sensing and boundary data \\
\hline
\end{tabular}
\end{table}

\paragraph{Mechanisation status.}
The three case-study families instantiate the same factorisation pattern.
The consensus case study mechanises both repair paths: the refinement
path restores admissibility for \(\texttt{quorum\_safe}\) by enriching
the observation (\texttt{ConsensusAdmissibility.v}), and the restriction
path yields an admissible local-quorum-belief predicate at the coarse
observational grain (\texttt{ConsensusRestriction.v}). The GPU case study
mechanises both repair paths: refinement by enriched global execution
observation (\texttt{GPUAdmissibility.v}), and restriction to the subtype
of data-race-free executions (\texttt{GPURestriction.v}). The
physical-systems case study mechanises both paths: a restricted predicate
is admissible at the coarse sensor grain, and an enriched sensor observation
restores admissibility for the stronger predicate. The development
mechanises predicate restriction in all three domains and observational
refinement in all three.

% -----------------------------------------------------------------------
\subsection{Consensus Under Partition}

The consensus case study concerns quorum safety under network partition. The
concrete state records enough information to determine whether the relevant
quorum-safety predicate holds. A coarse observation, however, may expose only
local status information. Such an observation does not determine global quorum
safety. Two concrete states may agree on the local observation while disagreeing
on whether a quorum is safely established.

This failure is not a defect in the proof script. It is the phenomenon the
framework is intended to detect. A uniformity principle for reachable alive
nodes fails in the partition model: different nodes may have different
reachable alive counts. The admissibility proof cannot be closed by that
route.

The mechanised consensus case study proves both repair paths.

\textit{Refinement.} The observation record \texttt{ConsensusObs} contains the
field \texttt{co\_quorum\_safe}, computed by \texttt{consensus\_obs} from the
underlying consensus state. The proof of admissibility is then immediate:
equality of observations entails equality of the quorum-safety component by
projection. The safety predicate becomes admissible because the observation
carries the safety-relevant distinction directly.

\textit{Restriction.} Rather than asserting global quorum safety, one can
instead assert only what the local observation supports: the node's own quorum
belief, \texttt{local\_quorum\_belief~n}, a projection of \texttt{lv\_quorum}
from the local view record. This predicate is admissible relative to
\texttt{local\_node\_obs~n} because it is a projection of the observation.
The proof requires no axiom. The theorem \texttt{restriction\_is\_genuinely\_weaker}
confirms that local belief and global quorum safety can disagree, so this is a
genuine weakening and not a reformulation.

The theorem \texttt{both\_repairs\_are\_available} states the formal duality
in a single proposition: both \texttt{local\_quorum\_belief~n} and
\texttt{quorum\_safe} are admissible, the former relative to the coarse local
observation, the latter relative to the enriched consensus observation.

% -----------------------------------------------------------------------
\subsection{GPU-Style Memory Observations}

The GPU-style case study concerns the gap between value observations and memory
consistency. A coarse observation may record only the values read and written by
threads. The concrete execution state, however, contains additional ordering and
visibility structure: buffering, cache state, scope, and synchronisation
metadata. A predicate such as sequential consistency of the concurrent execution
is not, in general, determined by the value trace alone.

The admissibility failure has the same form as in the consensus case. Two
executions may have the same observed value trace for a given thread while
disagreeing on whether the execution is sequentially consistent. The observation
map has collapsed a distinction on which the safety predicate depends. The
predicate therefore fails to factor through the thread-local observation.

The GPU case study mechanises both repair paths. Refinement restores
admissibility by enriching the observation with global execution structure.
Restriction keeps the thread-local observation fixed but restricts the domain
to data-race-free executions. On that restricted domain, the
sequential-consistency predicate is stable across observation fibres,
conditional on the DRF-SC bridge stated in \texttt{GPURestriction.v}.

% -----------------------------------------------------------------------
\subsection{Physical Sensing and Boundary Conditions}

The physical-systems case study concerns safety predicates over continuous or
boundary-sensitive states. A controller or monitor rarely observes the full
physical trajectory. It receives sampled, discretised, or otherwise compressed
sensor data. A predicate may require that a variable remain within a safe region
throughout an interval, while the observation records only values at selected
sampling points or coarse boundary summaries.

The resulting admissibility failure is again a failure of factorisation. Two
physical trajectories may induce the same sampled observation while disagreeing
on the safety predicate. One trajectory may remain inside the safe envelope,
while another may violate the envelope between samples or across an unresolved
boundary. The sampled observation is therefore insufficient to determine the
predicate.

This case study mechanises both recovery paths. The coarse sensor observation is
proved inadmissible for the original safety predicate. A restricted predicate
meaningful at the coarse observational grain is then shown to be admissible
relative to that observation. Finally, an enriched sensor observation that adds
the missing physical data restores admissibility for the stronger predicate. The
physical-systems case study therefore demonstrates the full design choice: one
may weaken the guarantee to match the observation, or strengthen the observation
to match the guarantee.

% -----------------------------------------------------------------------
\subsection{Common Structure}

The three case-study families differ in engineering content but not in logical
shape. Consensus hides partition and quorum structure. GPU-style memory hides
ordering and visibility structure. Physical sensing hides inter-sample or
boundary behaviour. In each case, the initial observation map is too coarse for
the safety predicate. The predicate fails to factor because the observation
identifies concrete states that the predicate separates.

This is the common admissibility defect. It is independent of whether the hidden
structure is distributed, concurrent, or physical. The repair is also common:
either choose a predicate that is meaningful at the given observational grain, or
refine the observation until the predicate becomes invariant on observation
fibres. The case studies therefore support the main methodological claim of the
paper. Admissibility is not merely a property of a predicate or a system. It is a
property of the relation between a predicate and an observation interface.

% at the point where the case-study section begins.

\section{Case Studies: Refine or Reject}

The development contains three case-study families: consensus under network
partition, GPU-style relaxed memory, and cyber-physical sensor models. These
examples are deliberately drawn from different engineering domains. Their common
structure is the point. In each case, a safety predicate is naturally formulated
over a concrete state space~\(X\), while the available observation map
\(M : X \to O\) exposes only a coarser interface. The case study then asks
whether the predicate factors through that observation.

The answer is negative for the initial coarse observation in each family. The
coarse observation identifies concrete states that disagree on the relevant
safety predicate. The admissibility failure is therefore structural rather than
accidental. Recovery takes one of two forms: replace the predicate by one that
is meaningful at the observational grain, or refine the observation so that the
safety-relevant distinction is no longer collapsed.

\begin{table}[t]
\centering
\caption{Case-study structure: factorisation failure and admissibility recovery.}
\label{tab:case-study-factorisation}
\small
\begin{tabular}{p{1.8cm}p{2.8cm}p{3.5cm}p{2.9cm}}
\hline
Case study & Hidden distinction & Coarse failure & Recovery \\
\hline
Consensus
  & Partition and quorum structure
  & Local observations do not determine global quorum safety
  & Restrict to local quorum belief; or enrich with quorum-safety data \\[4pt]
GPU memory
  & Buffer, cache, and ordering structure
  & Value traces do not determine consistency behaviour
  & Restrict to DRF executions; or enrich with global ordering metadata \\[4pt]
Physical systems
  & Inter-sample and boundary behaviour
  & Sampled values do not determine continuous safety
  & Use robust observational margins; or enrich sensing and boundary data \\
\hline
\end{tabular}
\end{table}

\paragraph{Mechanisation status.}
The three case-study families instantiate the same factorisation pattern.
The consensus case study mechanises both repair paths: the refinement
path restores admissibility for \(\texttt{quorum\_safe}\) by enriching
the observation (\texttt{ConsensusAdmissibility.v}), and the restriction
path yields an admissible local-quorum-belief predicate at the coarse
observational grain (\texttt{ConsensusRestriction.v}). The GPU case study
mechanises both repair paths: refinement by enriched global execution
observation (\texttt{GPUAdmissibility.v}), and restriction to the subtype
of data-race-free executions (\texttt{GPURestriction.v}). The
physical-systems case study mechanises both paths: a restricted predicate
is admissible at the coarse sensor grain, and an enriched sensor observation
restores admissibility for the stronger predicate. The development
mechanises predicate restriction in all three domains and observational
refinement in all three.

% -----------------------------------------------------------------------
\subsection{Consensus Under Partition}

The consensus case study concerns quorum safety under network partition. The
concrete state records enough information to determine whether the relevant
quorum-safety predicate holds. A coarse observation, however, may expose only
local status information. Such an observation does not determine global quorum
safety. Two concrete states may agree on the local observation while disagreeing
on whether a quorum is safely established.

This failure is not a defect in the proof script. It is the phenomenon the
framework is intended to detect. A uniformity principle for reachable alive
nodes fails in the partition model: different nodes may have different
reachable alive counts. The admissibility proof cannot be closed by that
route.

The mechanised consensus case study proves both repair paths.

\textit{Refinement.} The observation record \texttt{ConsensusObs} contains the
field \texttt{co\_quorum\_safe}, computed by \texttt{consensus\_obs} from the
underlying consensus state. The proof of admissibility is then immediate:
equality of observations entails equality of the quorum-safety component by
projection. The safety predicate becomes admissible because the observation
carries the safety-relevant distinction directly.

\textit{Restriction.} Rather than asserting global quorum safety, one can
instead assert only what the local observation supports: the node's own quorum
belief, \texttt{local\_quorum\_belief~n}, a projection of \texttt{lv\_quorum}
from the local view record. This predicate is admissible relative to
\texttt{local\_node\_obs~n} because it is a projection of the observation.
The proof requires no axiom. The theorem \texttt{restriction\_is\_genuinely\_weaker}
confirms that local belief and global quorum safety can disagree, so this is a
genuine weakening and not a reformulation.

The theorem \texttt{both\_repairs\_are\_available} states the formal duality
in a single proposition: both \texttt{local\_quorum\_belief~n} and
\texttt{quorum\_safe} are admissible, the former relative to the coarse local
observation, the latter relative to the enriched consensus observation.

% -----------------------------------------------------------------------
\subsection{GPU-Style Memory Observations}

The GPU-style case study concerns the gap between value observations and memory
consistency. A coarse observation may record only the values read and written by
threads. The concrete execution state, however, contains additional ordering and
visibility structure: buffering, cache state, scope, and synchronisation
metadata. A predicate such as sequential consistency of the concurrent execution
is not, in general, determined by the value trace alone.

The admissibility failure has the same form as in the consensus case. Two
executions may have the same observed value trace for a given thread while
disagreeing on whether the execution is sequentially consistent. The observation
map has collapsed a distinction on which the safety predicate depends. The
predicate therefore fails to factor through the thread-local observation.

The GPU case study mechanises both repair paths. Refinement restores
admissibility by enriching the observation with global execution structure.
Restriction keeps the thread-local observation fixed but restricts the domain
to data-race-free executions. On that restricted domain, the
sequential-consistency predicate is stable across observation fibres,
conditional on the DRF-SC bridge stated in \texttt{GPURestriction.v}.

% -----------------------------------------------------------------------
\subsection{Physical Sensing and Boundary Conditions}

The physical-systems case study concerns safety predicates over continuous or
boundary-sensitive states. A controller or monitor rarely observes the full
physical trajectory. It receives sampled, discretised, or otherwise compressed
sensor data. A predicate may require that a variable remain within a safe region
throughout an interval, while the observation records only values at selected
sampling points or coarse boundary summaries.

The resulting admissibility failure is again a failure of factorisation. Two
physical trajectories may induce the same sampled observation while disagreeing
on the safety predicate. One trajectory may remain inside the safe envelope,
while another may violate the envelope between samples or across an unresolved
boundary. The sampled observation is therefore insufficient to determine the
predicate.

This case study mechanises both recovery paths. The coarse sensor observation is
proved inadmissible for the original safety predicate. A restricted predicate
meaningful at the coarse observational grain is then shown to be admissible
relative to that observation. Finally, an enriched sensor observation that adds
the missing physical data restores admissibility for the stronger predicate. The
physical-systems case study therefore demonstrates the full design choice: one
may weaken the guarantee to match the observation, or strengthen the observation
to match the guarantee.

% -----------------------------------------------------------------------
\subsection{Common Structure}

The three case-study families differ in engineering content but not in logical
shape. Consensus hides partition and quorum structure. GPU-style memory hides
ordering and visibility structure. Physical sensing hides inter-sample or
boundary behaviour. In each case, the initial observation map is too coarse for
the safety predicate. The predicate fails to factor because the observation
identifies concrete states that the predicate separates.

This is the common admissibility defect. It is independent of whether the hidden
structure is distributed, concurrent, or physical. The repair is also common:
either choose a predicate that is meaningful at the given observational grain, or
refine the observation until the predicate becomes invariant on observation
fibres. The case studies therefore support the main methodological claim of the
paper. Admissibility is not merely a property of a predicate or a system. It is a
property of the relation between a predicate and an observation interface.


\subsection{Comparing the burden across domains}
\label{subsec:domain-burden}

The case studies share a formal pattern. They do not share a mathematical
substrate. In the consensus example, the relevant collapse is induced by local
views under partition. In the GPU example, the relevant collapse is induced by
thread-local observation of an execution whose global consistency status may
differ. In the physical-systems example, the collapse is mediated by sampling,
boundary behaviour, and mereotopological assumptions about regions.

The physical case is therefore the heaviest case, not because the criterion
changes, but because the observation boundary is richer. Spatial connection,
proper containment, inter-sample behaviour, and robust margins must be stated
before the factorisation question is even well posed. This is why the physical
case has explicit RCC-style assumptions and application-level coherence
constraints. The discrete cases need less geometric structure: they are simpler tests
of the same criterion.

The case studies exercise both repair forms in all three domains. This
symmetry is a mechanisation fact, not an engineering equivalence. In consensus,
restriction becomes a local-quorum-belief predicate. In GPU-style memory,
restriction becomes a data-race-free execution subdomain. In the physical case,
restriction becomes a robust observational margin. The formal pattern is common.
The cost of the repair is domain-specific.

This also explains the different burden of the restriction proofs. In the
consensus case, restriction changes the claim from global quorum safety to local
quorum belief. In the GPU case, restriction changes the domain from arbitrary
executions to data-race-free executions, with the DRF-SC bridge stated as a
proof obligation for the toy model. In the physical case, restriction is
naturally expressed by margins and boundary conditions. The same repair form is
present in all three cases. It is not the same engineering act in all three
cases.

The non-degenerate-threshold parameterisation clarifies this asymmetry.  The
physical case does not require special pleading.  It requires explicit
parameterisation of the arithmetic and boundary conditions that make its
observation space meaningful.  The GPU case carries an analogous burden
through the DRF-SC bridge.  The consensus case carries it through the
relation between local belief and global quorum structure.  The common
criterion is unchanged.  What changes is the cost of making each observation
boundary precise.

\subsection{Structure of the observation-induced equivalences}
\label{subsec:equivalence-structure}

The three case studies induce observation-equivalences with different
mathematical structures, despite sharing the same logical form
\(x \sim_M y \;\Leftrightarrow\; M(x)=M(y)\). Making this contrast explicit
clarifies why the physical case requires heavier axiomatisation and why
restriction is more natural there.

\medskip\noindent
\textbf{Consensus.} The equivalence \(ns_1 \sim_{M_n} ns_2\), where
\(M_n = \texttt{local\_node\_obs}\;n\), holds when both network states produce
the same \texttt{LocalView} record for node~\(n\). That record contains only a
message list filtered by the partition relation and a local quorum boolean. Two
states fall in the same fibre when node~\(n\) cannot observe the global
reachability information that separates them. The equivalence is combinatorial:
it depends entirely on which messages survive a single Boolean filter.

\medskip\noindent
\textbf{GPU memory.} The equivalence \(e_1 \sim_{T_t} e_2\), where
\(T_t = \texttt{thread\_local\_obs}\;t\), holds when both executions produce
the same sequence of memory events for thread~\(t\). The equivalence is
sequential: it collapses all inter-thread ordering, buffering, and
synchronisation metadata. Two executions fall in the same fibre when the
thread-local value trace is identical but the global consistency status differs.
No axioms beyond those for list equality are required to state this equivalence.

\medskip\noindent
\textbf{Physical sensing.} The equivalence \(s_1 \sim_{\sigma} s_2\), where
\(\sigma = \texttt{sensor\_obs}\), holds when both physical states share the
same boundary pressure and flux readings, regardless of temperature. This is
the simplest fibre condition to write down, but it is not the cheapest to reason
about. Constructing the inadmissibility witness requires that
\(\texttt{max\_temp} - 1 < \texttt{max\_temp}\), which is only valid under the
positivity axiom. Stating what a safe region \emph{is}, before asking whether
it is observable, requires the boundary and containment structure provided by
the RCC-style parameters and axioms. The equivalence is geometrically
grounded in a way the discrete equivalences are not.

\medskip
The contrast is this. In the discrete cases the observation-induced equivalence
is structurally transparent: list filtering and record projection. In the
physical case the equivalence requires background geometric conditions to be
non-degenerate. This is not a weakness of the physical model. It reflects the
genuine cost of stating a spatial safety predicate precisely enough to ask the
factorisation question at all.

% -----------------------------------------------------------------------
\section{Related Work}
\label{sec:related}
% -----------------------------------------------------------------------

\subsection{Abstraction and abstract interpretation}

The nearest technical setting is abstract interpretation. Cousot and Cousot
define program analysis by interpreting concrete computations in an abstract
domain, with soundness controlled by the relation between concrete and abstract
semantics~\cite{cousot1977abstract}. The present paper uses only the following
fragment of that setting: a predicate must be constant on the observational
fibres if it is to be asserted through the observation.

Counterexample-guided abstraction refinement gives a standard method for
repairing abstractions that admit spurious counterexamples~\cite{clarke2000cegar,clarke2003cegar}.
The criterion here is prior to refinement search: before asking how to refine,
one asks whether the asserted predicate is even invariant on the current
observational fibres. CEGAR asks how to refine after an abstract counterexample.
Structural admissibility asks whether the predicate was observable through the
interface in the first place.

\subsection{Runtime verification and monitorability}

Runtime verification studies what can be judged from finite or partial runtime
observations. Bauer, Leucker, and Schallhart give a three-valued account for LTL
and TLTL monitoring, including inconclusive outcomes over partial
traces~\cite{bauer2011runtime}. This paper does not replace monitorability theory.
It states a basic obstruction that any such theory must respect: a monitor cannot
distinguish states that its observation interface identifies.

Functional safety and risk-management practice already distinguish between
changing the operating claim and changing the instrumentation. Standards such
as ISO~31000 and IEC~61508 are not factorisation theories, but they make clear
that risk treatment, safety functions, monitoring, and operating conditions must
be stated together~\cite{iso31000,iec61508}. This paper supplies a small formal
criterion for one version of that practical discipline.

\subsection{Consensus and weak memory}

The consensus case study is intentionally small. It should be read against the
larger literature on distributed agreement, including Lamport's presentation of
Paxos~\cite{lamport2001paxos}. The point is not to mechanise Paxos. The point is
to show how a local observation can collapse states that differ with respect to
a global safety predicate.

The GPU-style case is similarly modest. It is motivated by the general problem
that weak-memory behaviours depend on global consistency constraints not
visible in a local thread trace. Work on x86-TSO and weak memory models gives
the broader context for this issue~\cite{owens2009x86tso,alglave2014herding}.
Here again, the paper does not claim to mechanise a production memory model. It
uses the setting to expose factorisation failure.

\subsection{Type-theoretic and constructive verification}

The mechanisation is carried out in Rocq, formerly Coq. Proof assistants make
assumptions, definitions, and proof obligations explicit. That explicitness is
part of the paper's method. A hidden observational assumption is not harmless
because it permits a theorem to compile. It changes the observation through
which the theorem is being claimed~\cite{coq2024}.

\subsection{Mereotopology and region-based spatial reasoning}

The physical-systems case uses region-based reasoning in the spirit of the
Region Connection Calculus. Randell, Cui, and Cohn's RCC framework treats
spatial reasoning in terms of regions and connection rather than point-set
coordinates~\cite{randell1992rcc}. The present paper uses only a small part of
this tradition. Its role is to make boundary-sensitive observation explicit.

% -----------------------------------------------------------------------
\section{Limitations}
\label{sec:limitations}
% -----------------------------------------------------------------------

The case studies are small by design. Their role is separation, not coverage:
each isolates a concrete observation map, a predicate, a collapse witness, and
a repair. No result depends on treating these models as complete industrial
specifications.

Both repair forms are mechanised in all three domains. This is a coverage
claim about the case studies, not a claim that the repairs have equal strength,
equal cost, or equal naturalness in each domain. Consensus formalises
restriction (\texttt{ConsensusRestriction.v}) and refinement
(\texttt{ConsensusAdmissibility.v}). GPU memory formalises restriction
(\texttt{GPURestriction.v}) and refinement (\texttt{GPUAdmissibility.v}).
The physical-systems case formalises both robust restriction and refinement.

The GPU restriction result should not be read as a full mechanisation of a
production GPU memory model or as an independent proof of a DRF-SC theorem.
It proves the admissibility repair obtained once the execution class is
restricted to data-race-free executions and the DRF-SC bridge is supplied for
the toy model. This is the intended level of the case study.

The RCC axioms and application-level coherence constraints are explicit
mathematical assumptions. They are not derived from a minimal connection theory.
The paper gives a structural criterion, not an automated refinement algorithm
and not a decision procedure for admissibility.

Finally, warrant debt is not proposed as a quantitative risk metric here.
The development gives warrant debt a qualitative profile: debt kind, repair pressure,
and repair cost band.  This is enough to distinguish trace debt from boundary
debt, or threshold debt from runtime debt.  It is not enough to compute
expected harm, probability of failure, or economic cost.  Those quantitative
questions remain future work.

% -----------------------------------------------------------------------
\section{Discussion}
\label{sec:discussion}
% -----------------------------------------------------------------------

\subsection{Observation interfaces as verification boundaries}

A safety property is a property of the relation between the system, the
predicate, and the observation interface. If the observation cannot distinguish
two states, no claim that separates those states is justified through that
observation.

\subsection{Refinement as factorisation repair}

Refinement is often described as the removal of spurious behaviour from an
abstraction. Here it repairs a failed factorisation: it changes the observation
boundary so that a predicate invisible at the coarse interface becomes visible
at the refined one.

This way of speaking also clarifies why refinement is not always the right
repair. Sometimes the better engineering decision is to restrict the claim.
Operating envelopes, robust margins, lockout conditions, and restricted modes
are not second-class repairs. They are ordinary safety devices. What matters is
that the restricted claim actually factors through the observation being used.

\subsection{Warrant debt}

Warrant debt is the operational residue of claiming more than the observation
can support.  The formal object is the collapse witness: same observation,
different predicate value.  The term is useful only if it prevents
equivocation.  If the missing distinction is irrelevant to the predicate,
there is no debt.  If the predicate requires the missing distinction, then
calling the discrepancy noise is unsafe.

The development treats warrant debt as a qualitative classification,
not merely as a name.  A debt profile records the kind of collapsed
distinction, the repair pressure, and the expected repair cost band.  Quorum
debt, trace debt, boundary debt, threshold debt, and runtime debt all
instantiate the same factorisation failure, but they do not impose the same
engineering burden.

This classification also clarifies what it would mean to carry warrant debt.
A project may decide that full observational refinement is too costly.  That
does not discharge the debt.  It changes the safety case.  The stronger claim
must be rejected or restricted, and the remaining debt must be recorded as a
limitation of the supported claim.  In this sense, acceptable warrant debt is
not hidden debt.  It is named, bounded, and attached to a weaker assertion.

The proof obligation remains unchanged: show factorisation, restrict the
claim, reject the claim, or refine the observation.

\paragraph{Future work.}
Future work has three natural directions. The first is quantitative: to refine
the qualitative warrant-debt profile into a numerical account of severity, for
example by metric distance, information-theoretic loss, or expected repair cost.
The second is mechanistic: to replace the DRF-SC bridge in the GPU case with a
fuller proof inside a richer memory model. The third is scaling: to move from
deliberately small case studies to larger assurance models without changing the
factorisation criterion.

% -----------------------------------------------------------------------
\section{Conclusion}
\label{sec:conclusion}
% -----------------------------------------------------------------------

The paper proves a limited claim. Its consequence is not optional. A safety
predicate can be asserted through an observation only if it factors through
that observation. If two concrete states have the same observation but different
predicate values, the observation does not support the claim. Post-processing
cannot repair that loss.

This situation occurs in ordinary safety work. Engineers restrict operating
envelopes. They add sensors. They change thresholds. They refine logs. They
reject claims that the evidence cannot support. The paper does not replace
those practices. It gives one formal reason why they matter.

The implemented lesson is therefore small but strict.  A failed
factorisation is not an invitation to tune a post-processor.  It is a repair
obligation.  The obligation may be discharged by restriction, rejection, or
refinement.  If it is not discharged, it must be carried explicitly as profiled
warrant debt.

The warning is simple. An admissibility failure may look like noise. It may look
like harmless underspecification. It may look like a modelling convention. But
if the safety predicate depends on a distinction the observation has collapsed,
the problem is structural. The response is not a stronger post-processor. It
is restriction, rejection, or refinement.

% -----------------------------------------------------------------------
\appendix

\section{Assumptions}
\label{app:assumptions}
% -----------------------------------------------------------------------

The development contains no admitted lemmas, 41 compiled Rocq source files, and
eight named axioms. Every axiom is documented in \texttt{ASSUMPTIONS.md} in the
accompanying repository. The audit commands that verify this status are:

\begin{verbatim}
find . -name '*.v'  | wc -l   # 41
find . -name '*.vo' | wc -l   # 41
grep -RniE '\bAdmitted\b' . --include='*.v'   # no output
grep -RnE  '\bAxiom\b'    . --include='*.v'   # 8 matches
\end{verbatim}

\noindent Both grep commands above also return no output when run against
\texttt{GPURestriction.v} directly, confirming that the restriction repair
introduces neither \texttt{Axiom} declarations nor \texttt{Admitted} proofs.

% -----------------------------------------------------------------------
\section{Rocq File Map}
\label{app:filemap}
% -----------------------------------------------------------------------

\begin{table}[h]
\centering
\small
\begin{tabular}{p{2.1cm}p{4.0cm}p{4.5cm}}
\toprule
Component & Files & Role \\
\midrule
Foundation & \texttt{Foundation/*.v} & Observation, factorisation, mereotopology \\
Admissibility & \texttt{Admissibility/*.v} & Core admissibility theory \\
Runtime & \texttt{Runtime/*.v} & Monitor states, fatigue, certificates \\
Automation & \texttt{Automation/*.v} & Custom proof tactics \\
Consensus & \texttt{CaseStudies/\allowbreak{}Consensus/*.v} & Quorum-safety admissibility, local-quorum restriction \\
GPU & \texttt{CaseStudies/GPU/*.v} & Thread-local inadmissibility, global refinement, DRF restriction \\
Physical & \texttt{CaseStudies/\allowbreak{}PhysicalSystems/*.v} & Sensor inadmissibility, restriction, refinement \\
Extraction & \texttt{Extraction/*.v} & Verified OCaml extraction bridge \\
\bottomrule
\end{tabular}
\caption{Rocq source file map.}
\label{tab:filemap}
\end{table}

\FloatBarrier

% -----------------------------------------------------------------------
% Bibliography
% -----------------------------------------------------------------------
\begin{thebibliography}{99}

\bibitem{alglave2014herding}
Jade Alglave, Luc Maranget, and Michael Tautschnig.
\newblock Herding cats: modelling, simulation, testing, and data mining for weak memory.
\newblock \emph{ACM Transactions on Programming Languages and Systems}, 36(2), 2014.

\bibitem{bauer2011runtime}
Andreas Bauer, Martin Leucker, and Christian Schallhart.
\newblock Runtime verification for LTL and TLTL.
\newblock \emph{ACM Transactions on Software Engineering and Methodology}, 20(4):14:1--14:64, 2011.

\bibitem{clarke2000cegar}
Edmund M. Clarke, Orna Grumberg, Somesh Jha, Yuan Lu, and Helmut Veith.
\newblock Counterexample-guided abstraction refinement.
\newblock In \emph{Computer Aided Verification}, LNCS 1855, pages 154--169. Springer, 2000.

\bibitem{clarke2003cegar}
Edmund M. Clarke, Orna Grumberg, Somesh Jha, Yuan Lu, and Helmut Veith.
\newblock Counterexample-guided abstraction refinement for symbolic model checking.
\newblock \emph{Journal of the ACM}, 50(5):752--794, 2003.

\bibitem{cousot1977abstract}
Patrick Cousot and Radhia Cousot.
\newblock Abstract interpretation: a unified lattice model for static analysis of programs by construction or approximation of fixpoints.
\newblock In \emph{Proceedings of the 4th ACM Symposium on Principles of Programming Languages}, pages 238--252, 1977.

\bibitem{iec61508}
International Electrotechnical Commission.
\newblock \emph{IEC 61508: Functional safety of electrical/electronic/programmable electronic safety-related systems}.
\newblock IEC, current multi-part standard.

\bibitem{iso31000}
International Organization for Standardization.
\newblock \emph{ISO 31000:2018 Risk management, Guidelines}.
\newblock ISO, 2018.

\bibitem{lamport2001paxos}
Leslie Lamport.
\newblock Paxos made simple.
\newblock \emph{ACM SIGACT News}, 32(4):51--58, 2001.

\bibitem{owens2009x86tso}
Scott Owens, Susmit Sarkar, and Peter Sewell.
\newblock A better x86 memory model: x86-TSO.
\newblock In \emph{Theorem Proving in Higher Order Logics}, LNCS 5674, pages 391--407. Springer, 2009.

\bibitem{randell1992rcc}
David A. Randell, Zhan Cui, and Anthony G. Cohn.
\newblock A spatial logic based on regions and connection.
\newblock In \emph{Proceedings of the 3rd International Conference on Principles of Knowledge Representation and Reasoning}, pages 165--176, 1992.

\bibitem{coq2024}
The Coq Development Team.
\newblock \emph{The Coq Proof Assistant}.
\newblock Zenodo, 2024.

\end{thebibliography}

\end{document}
